\section{Aufbau der Arbeit}

Die Arbeit gliedert sich in fünf Kapitel. Nach der Einleitung werden in Kapitel 2 die literaturbasierten Ausgangspunkte dargestellt. 
Neben der Definition und den Besonderheiten von KMU werden dabei der Digitalisierungsdruck unter begrenzten Ressourcen, die 
Rolle von ERP-Systemen als Enabler sowie das ERP-Paradoxon betrachtet. Darauf aufbauend werden zentrale Voraussetzungen für 
die Wirksamkeit von ERP-Systemen in KMU herausgearbeitet.

Kapitel 3 überführt diese Voraussetzungen in das konzeptionelle ERP-KMU-Fit-Modell. Dazu werden die Prozess-, Daten-, Ressourcen-, 
Kompetenz- und Strategie-Dimension einzeln erläutert. Anschließend wird dargestellt, wie diese Dimensionen zusammenwirken und die 
Implementierungsbereitschaft eines KMU beeinflussen.

Kapitel 4 behandelt die Grenzen und Limitationen des entwickelten Modells. 

Kapitel 5 fasst die zentralen Erkenntnisse zusammen und beantwortet die Forschungsfrage.